% Agent /root/monodromy; complete argument, 19 November 2026.
% Independently audited in full by /root, including the opened classification source.
\subsection{Plane curves in odd characteristic: the full group in every degree}
\label{sec:odd-plane-all-degrees}
\status{Proved; independently audited by /root, including the primary classification statement.}

Throughout this section $p$ is odd, $k$ is an algebraic closure of
$\mathbf F_p$, $d\geq4$, and $g=(d-1)(d-2)/2$.  Let $W_d$ be the
Hasse--Witt invertibility open in the projective space of degree-$d$
plane equations, including singular equations, and let $U_d\subset W_d$
be its smooth open.  The vector bundle $R^1\pi_*\mathcal O$ has constant
rank $g$ on the entire projective parameter space.  Its Lang fixed
sheaf $V$ is finite locally constant of rank $g$ on $W_d$.
The map $\pi_1(U_d)\to\pi_1(W_d)$ is surjective.  Thus computing its
image on $W_d$ computes the desired geometric image on $U_d$.
Choose as base point a union of $d$ lines with no triple intersection.
All groups below act on this same fibre, after choosing paths in the
specified connected parameter spaces.  This is permissible because
every boundary family used below is contained in $W_d$.

\subsubsection{The line-arrangement representation}

Put $P=\mathbf F_p^d$, with its permutation basis $e_1,\ldots,e_d$,
$u=e_1+\cdots+e_d$, and $W=P/\langle u\rangle$.  The graph of the
line arrangement is the complete graph $K_d$.  Its first cochain space
is $\bigwedge^2P$, with $e_i\wedge e_j$ representing an oriented edge.
The coboundaries are $u\wedge P$.  Consequently
\begin{equation}\label{eq:odd-arrangement-module}
 V=H^1(K_d,\mathbf F_p)
   =\bigwedge^2P/(u\wedge P)=\bigwedge^2W.
\end{equation}
This is cohomology, rather than the dual cycle space.  The coherent
normalization sequence gives the same identification after tensoring
with $k$, with bijective Frobenius on all the graph terms.

The space of ordered general-position lines is irreducible; its
quotient by the freely acting group $S_d$ is the unordered arrangement
space.  Its ordered cover is connected, so its permutation monodromy
is $S_d$.  Over the ordered cover the graph local system is constant.
Thus the image $G\subseteq\mathrm{GL}(V)$ contains exactly the
permutation action in \eqref{eq:odd-arrangement-module}.

Let $U=\{(x_i)\in P:\sum x_i=0\}$, let $D$ be its image in $W$, and
put $A=\bigwedge^2D\subseteq\bigwedge^2W$.

\begin{lemma}\label{lem:odd-wedge-submodules}
Every nonzero $S_d$-submodule of $\bigwedge^2W$ contains $A$.
If $p\nmid d$, then $A=V$ and $V$ is simple.  If $p\mid d$, then
$A$ is the unique proper nonzero submodule, and $V/A\simeq D$ is
simple of dimension $d-2$.
\end{lemma}

\begin{proof}
For four distinct indices $i,j,k,l$, put $r=e_i-e_j$ and $s=e_k-e_l$.
The group-algebra operator
\[
 (1-(ij))(1-(kl))
\]
on $\bigwedge^2W$ has image $\langle r\wedge s\rangle$ and is the
rank-one map obtained by evaluating the functional
$(x_i-x_j)\wedge(x_k-x_l)$ and multiplying by $r\wedge s$.
The latter vector is nonzero: the line $\langle u\rangle$ does not
meet $\langle r,s\rangle$, including when $d=4$, since $p\ne2$.

The vectors $r\wedge s$ for disjoint pairs span $\bigwedge^2U$ before
passage to $W$, and hence span $A$ afterwards.  Explicitly, on four
indices take $f_i=e_i-e_l$, $f_j=e_j-e_l$, $f_k=e_k-e_l$.  The three
perfect matchings give the columns
\[
 \begin{pmatrix}0&1&1\\1&0&-1\\-1&1&0\end{pmatrix}
\]
in the basis $f_i\wedge f_j,f_i\wedge f_k,f_j\wedge f_k$; their
determinant is $2$.  Every pair of difference-basis vectors belongs to
such a four-index set.  The same argument in the dual permutation
space shows that the corresponding functionals span
$\bigwedge^2W^*$, because $W^*$ is the sum-zero subspace of $P^*$.
For every nonzero vector of $V$, therefore, one of these rank-one
operators has nonzero value.  Its entire $S_d$-orbit spans $A$.
This proves the first assertion and also proves that $A$ is simple.

If $p\nmid d$, then $U\to W$ is an isomorphism.  If $p\mid d$,
$D=U/\langle u\rangle$ has dimension $d-2$, and the sum functional
$\sigma$ descends to $W$ with kernel $D$.  Contraction with $\sigma$
gives the exact sequence
\[
 0\longrightarrow\bigwedge^2D\longrightarrow\bigwedge^2W
 \xrightarrow{\ \iota_\sigma\ }D\longrightarrow0.
\]
The module $D$ is simple: a nonzero class has a representative whose
coordinates are not all equal; applying $1-(ij)$ for two unequal
coordinates produces a nonzero multiple of $e_i-e_j$, and the
translates of this difference span $D$.  Every proper nonzero
submodule of $V$ contains $A$ and has zero image in the simple
quotient unless it is all of $V$.  This proves uniqueness.
\end{proof}

\begin{proposition}\label{prop:odd-plane-irreducible}
The geometric monodromy group $G$ is irreducible on $V$.
\end{proposition}

\begin{proof}
Only the case $p\mid d$ needs another family.  Take a smooth ordinary
curve of degree $d-2$ and a smooth conic meeting it transversely in
$m=2d-4$ points.  The graph contribution is a sub-local-system
\[
 B=\mathbf F_p^m/\langle(1,\ldots,1)\rangle\subset V
\]
of rank $m-1=2d-5$.  Its node monodromy is the full symmetric group
$S_m$.  To verify this, fix the conic.  Restriction of degree-$(d-2)$
plane forms to it is surjective onto binary forms of degree $m$.
Over the squarefree binary forms, the ordered-root cover is connected.
Pulling it back gives an affine bundle after choosing a coefficient
chart (equivalently, an open of a projective bundle).  Imposing
smoothness and ordinarity of the plane curve removes an open subset
from this irreducible total space, and the open is nonempty.  The
ordered-root cover remains connected.  Thus the induced graph image
is $S_m$.

Here $p\nmid m$, since $p\mid d$ and $p$ is odd.  The deleted
permutation module $B$ is simple, by the same difference argument
used above.  If the only possible proper $G$-invariant subspace $A$
were invariant, the map $B\to V/A$ would be zero: its source is a
simple module of dimension $2d-5>d-2=\dim(V/A)$.  Hence $B\subset A$.

This inclusion can be tested at the common line arrangement.  Allow
the conic to become the union of the first two lines and the other
curve to become the remaining $d-2$ lines.  Keep the intersections
between the two subcurves transverse, disjoint from their internal
nodes.  Their ordinary coherent cohomology and the graph sequence
extend over an open containing both this point and the smooth
subcurve locus.  At this arrangement $B$ is the kernel of restriction
to the complete subgraph on vertices $3,\ldots,d$.  In particular the
edge cochain $e_1\wedge e_3$ represents an element of $B$.  But
\[
 \iota_\sigma(e_1\wedge e_3)=e_3-e_1\ne0\quad\hbox{in }D,
\]
so it does not belong to $A$.  This contradiction proves the claim.
No specialization through a drop of the Hasse--Witt rank has been
used.
\end{proof}

\subsubsection{Quartic boundaries and an elementary cocycle calculation}

\begin{lemma}\label{lem:sl3-natural-h1-odd}
For every odd prime $p$, the natural three-dimensional module $Q$
and its dual satisfy
$H^1(\mathrm{SL}_3(\mathbf F_p),Q)=
H^1(\mathrm{SL}_3(\mathbf F_p),Q^*)=0$.
\end{lemma}

\begin{proof}
Let $c$ be a cocycle.  The diagonal sign subgroup of determinant one
has order four, prime to $p$, so subtracting a coboundary makes $c$
zero on it.  Each $t_{ij}(a)=1+aE_{ij}$ commutes with the sign matrix
which is $-1$ in positions $i,j$ and $+1$ in the remaining position
$k$.  The cocycle relation then forces
$c(t_{ij}(a))\in\langle e_k\rangle$.
Write $x=t_{12}(a)$, $y=t_{23}(b)$, with $a,b\ne0$, and
$c(x)=Ae_3$, $c(y)=Be_1$.  Since $[x,y]=t_{13}(ab)$, the cocycle
identity gives
\[
 c([x,y])=-bA e_2-abA e_1.
\]
Its left side belongs to $\langle e_2\rangle$, hence $A=0$.
Permuting the three indices proves vanishing on every elementary
transvection.  They generate $\mathrm{SL}_3(\mathbf F_p)$, so the
normalized cocycle is zero.  Applying the inverse-transpose
automorphism proves the dual assertion.
\end{proof}

For any chosen four of the $d$ lines, fix the other $r=d-4$ lines and
vary an ordinary quartic, transverse to the fixed lines and avoiding
their pairwise intersections.  Include ordinary singular quartics
whenever the total degree-$d$ curve stays in $W_d$.  This connected
open contains the original arrangement.  Mayer--Vietoris for the
quartic and the fixed lines gives
\begin{equation}\label{eq:odd-quartic-boundary}
 0\longrightarrow T\longrightarrow V\longrightarrow Q
 \longrightarrow0,\qquad \dim Q=3,
\end{equation}
where $Q$ is the quartic Hasse--Witt local system.  Its image is
$\mathrm{GL}_3(\mathbf F_p)$ by the canonical quartic moduli argument
and Achter--Pries, Proposition~4.8, already checked above.  The image
on $T$ is solvable: its only moving graph edges are the four
intersections with each fixed line, so it factors through a subgroup
of $(S_4)^r$.  The nodes between the fixed lines are fixed.

Let $B$ be the image of this boundary family and let $D_T$ be its
image on $T$.  The image of $B$ in
$\mathrm{GL}(Q)\times D_T$ is a subdirect product.  Goursat's lemma,
or the elementary description of a subdirect product as equality in
a common quotient, shows that it contains
$\mathrm{SL}(Q)\times1$: every common quotient is solvable, whereas
$\mathrm{SL}_3(\mathbf F_p)$ is perfect.  Let $B_0$ be the inverse
image of $\mathrm{SL}(Q)\times1$.  It is normal in $B$, fixes $T$
pointwise, and maps onto $\mathrm{SL}(Q)$.
Its kernel is an additive submodule
\begin{equation}\label{eq:odd-quartic-kernel}
 K\subseteq\operatorname{Hom}(Q,T)=Q^*\otimes T.
\end{equation}
The natural dual is absolutely simple, with endomorphism ring
$\mathbf F_p$.  A submodule of this direct sum of copies of $Q^*$
therefore has the form $K=Q^*\otimes T'$ for a subspace $T'\subseteq T$.
The subspace $T'$ is $B$-invariant because $K$ is normal in $B$.
If $K\ne0$, it contains rank-one maps $w\otimes\phi$ and hence
transvections $1+w\otimes\phi$ of $V$.  If $K=0$, $B_0$ is a
complement to $\operatorname{Hom}(Q,T)$ in the indicated semidirect
product.  Lemma~\ref{lem:sl3-natural-h1-odd}, applied to each copy of
$Q^*$, makes that complement conjugate to the block-diagonal
$\mathrm{SL}(Q)$.  In this case too it contains transvections of $V$.
For $d=4$ the statement simply recovers the full quartic group.

\subsubsection{The normal subgroup generated by transvections}

\begin{lemma}\label{lem:transvection-clifford-blocks}
Let an irreducible finite group $G\subseteq\mathrm{GL}(V)$ contain a
transvection, and let $N\triangleleft G$ be the subgroup generated by
all transvections in $G$.  Then $V$ is a direct sum of pairwise
nonisomorphic simple $N$-modules, and $G$ permutes these summands
transitively.  Every transvection has its image line in exactly one
summand and fixes all the others pointwise.
\end{lemma}

\begin{proof}
The sum of the simple $N$-submodules, $\operatorname{soc}_N(V)$, is
nonzero and $G$-stable, since $N$ is normal.  Irreducibility of $G$
therefore makes it all of $V$.  The canonical isotypic summands are
permuted transitively by $G$.  Write any one of them as $S^{\oplus e}$.
For an element of $N$, the rank of its difference from the identity
on this summand is a multiple of $e$.  Some transvection has nonzero
action on some summand, and its total rank is one, so $e=1$ there;
transitivity gives $e=1$ for all summands.  Thus the canonical
isotypic summands are themselves simple.  Each transvection is
block diagonal because it belongs to $N$, and its rank-one
difference is nonzero on only one block.  This proves the assertion.
(The same argument with the finite division endomorphism ring of a
block shows it is absolutely simple, but that strengthening is not
needed below.)
\end{proof}

For a four-element subset $I\subset\{1,\ldots,d\}$, let
$U_I\subset P$ be its three-dimensional sum-zero permutation
subspace, supported on $I$.  Because $4\ne0$ in $\mathbf F_p$,
$U_I$ embeds in $W$.  Set
\[
 S_I=\bigwedge^2 U_I\subset V.
\]
The subgroup $S_I^{\rm perm}\simeq S_4$ permuting these four lines
lies in the corresponding quartic boundary group $B$.  Restricted
to it, the decomposition of the permutation space gives
\begin{equation}\label{eq:s4-quartic-decomposition}
 W\simeq U_I\oplus\mathbf F_p^{d-4},\qquad
 V\simeq S_I\oplus U_I^{\oplus(d-4)}
             \oplus\mathbf F_p^{\binom{d-4}{2}}.
\end{equation}
Both $U_I$ and $S_I$ are simple.  For the former use the elementary
deleted-permutation argument with $p\nmid4$; for the latter use
Lemma~\ref{lem:odd-wedge-submodules} with four indices.  They are
nonisomorphic: a transposition has trace $1$ on $U_I$ and trace $-1$
on $\bigwedge^2U_I$, distinct in every odd characteristic.  Neither
is trivial.  Thus \eqref{eq:s4-quartic-decomposition} is a semisimple
decomposition even when $p=3$, and $S_I$ occurs with multiplicity one.
The quotient $Q$ in \eqref{eq:odd-quartic-boundary} is isomorphic to
$S_I$ as an $S_4$-module, and $T$ contains no constituent of this type.

\begin{proposition}\label{prop:odd-transvection-normal-irreducible}
The normal subgroup $N\triangleleft G$ generated by all transvections
is irreducible on $V$.
\end{proposition}

\begin{proof}
Use the canonical $N$-block decomposition of
Lemma~\ref{lem:transvection-clifford-blocks}.  We first prove that for
each quartet $I$ its subspace $S_I$ lies entirely in one block.

If the kernel $K$ of its boundary group $B_0$ is zero, the cocycle
calculation makes $B_0$ a conjugate block-diagonal $\mathrm{SL}_3$.
Its commutator subspace $P_I=[B_0,V]$ is therefore three dimensional,
simple as a $B_0$-module, and maps isomorphically to $Q$.  Normality
of $B_0$ in $B$ makes $P_I$ stable under the quartet's $S_4$.
The unique-copy statement in \eqref{eq:s4-quartic-decomposition}
forces $P_I=S_I$.  This copy of $\mathrm{SL}_3$ is generated by
rank-one transvections of $V$, so $B_0\subseteq N$.  Consequently
each $N$-block is $B_0$-stable and
$[B_0,V]=\bigoplus_j[B_0,V_j]$.  Simplicity of $P_I$ puts it in one
$N$-block.

If $K\ne0$, write it as $Q^*\otimes T'$, with $T'\ne0$, and fix
one nonzero $\phi\in Q^*$.  For every nonzero $w\in T'$, the
transvection $1+w\otimes\phi$ is supported on a single $N$-block.
The functional here means the pullback of $\phi$ along $V\to Q$.
It is nonzero on that block and zero on every other block.  Since
the same nonzero functional is used for all $w$, that block must
be the same block $L$ for every $w$.  Thus $T'\subset L$.
Now fix one nonzero $w$ and vary $\phi$.  Every $\phi\in Q^*$
vanishes on every other $N$-block.  It follows that $L\to Q$ is
surjective.  Since $T'$ is stable under the quartet's $S_4$, this
$S_4$ must preserve $L$: its translate is another $N$-block sharing
the nonzero subspace $T'$.  As an $S_4$-submodule of the semisimple
module \eqref{eq:s4-quartic-decomposition}, $L$ is semisimple.
Its surjection onto $Q\simeq S_I$ forces it to contain the unique
copy $S_I$ in $V$.  This proves the claim in both cases.

When two quartets have three common indices, their subspaces $S_I$
intersect in the nonzero line $\bigwedge^2U_{I\cap J}$.
The graph of four-element subsets joined by such overlaps is
connected.  Hence all $S_I$ belong to the same $N$-block $L$.
Their span is $A=\bigwedge^2D$: the spanning disjoint-difference
wedges in Lemma~\ref{lem:odd-wedge-submodules} are supported on
quartets.  Thus $A\subset L$.

The subspace $A$ is $S_d$-stable.  Every translate of $L$ by $S_d$
is an $N$-block containing $A$, hence is $L$ itself.  Thus $S_d$
preserves $L$ and the sum $L'$ of the other blocks.  If $L'\ne0$,
Lemma~\ref{lem:odd-wedge-submodules} forces $A\subset L'$, which
contradicts $A\subset L$ and $L\cap L'=0$.  There are no other
blocks, and $N$ is irreducible.  The base case $d=4$ was already
settled by the quartic monodromy theorem.
\end{proof}

\subsubsection{Classification, exclusion of the alternating form, and determinant}

McLaughlin's transvection theorem, in the precise form of
Kantor--Seress (2001), Theorem~3.6, p.~22 \cite{KantorSeress2001},
states that for $n\geq3$
an irreducible subgroup of $\mathrm{SL}_n(\mathbf F_q)$ generated by
full transvection groups is $\mathrm{SL}_n(\mathbf F_q)$ or, when
$n$ is even, $\mathrm{Sp}_n(\mathbf F_q)$; the additional alternatives
in that theorem occur only when $q=2$.  Here $q=p$ is prime, and
the powers of $1+v\phi$ are precisely $1+a v\phi$, $a\in\mathbf F_p$.
Thus every generating transvection supplies its full transvection
group.  The theorem applies to $N$, which has determinant one and
is irreducible by Proposition~\ref{prop:odd-transvection-normal-irreducible}.
It leaves $N=\mathrm{SL}_g(\mathbf F_p)$ or $N=\mathrm{Sp}_g(\mathbf F_p)$.

The symplectic possibility is impossible.  If $N$ is symplectic,
normality puts $G$ in its group of similitudes: the invariant
alternating form is unique up to scalar.  For any quartet, the
multiplier on $B_0$ is trivial.  Indeed it is trivial on its
$p$-group kernel $K$ and on its perfect quotient
$\mathrm{SL}_3(\mathbf F_p)$.  Therefore $B_0$ preserves the
nondegenerate alternating form.  Since it fixes $T$ pointwise and
has fixed-point-free natural quotient $Q$, its fixed space is
$V^{B_0}=T$.  For a group preserving a nondegenerate bilinear form,
\[
 [B_0,V]=(V^{B_0})^\perp=T^\perp.
\]
The latter space has dimension three.  Its image in $Q$ is
$[B_0,Q]=Q$, so it maps isomorphically to $Q$ and meets $T$ trivially.
The alternating form is consequently nondegenerate on this
three-dimensional space $T^\perp$, a contradiction.  We have proved
\begin{equation}\label{eq:odd-plane-sl-contained}
 \mathrm{SL}_g(\mathbf F_p)\subseteq G.
\end{equation}

It remains to show that the geometric determinant is all of
$\mathbf F_p^\times$.  Take a Legendre elliptic cubic $E_t$, near a
supersingular parameter $t_0$, and add $d-3$ fixed lines chosen in
general position relative to $E_{t_0}$.  The Hasse polynomial of
the Legendre family has a simple root at $t_0\notin\{0,1\}$, as
proved in the elliptic calculation in
\S\ref{subsec:plane-residues}.  Over the strict local
trait at $t_0$, all graph contributions have bijective Frobenius
and determinant a unit.  The normalization sequence therefore
makes the total Hasse determinant a unit times the elliptic Hasse
invariant.  It has valuation exactly one at $t_0$.
On the punctured trait the total curve is ordinary, hence belongs
to $W_d$.  The determinant Lang torsor is
$z^{p-1}=h^{-1}$ by \eqref{eq:determinant-kummer}.  A valuation-one
$h$ gives tame inertia image of exact order $p-1$.  This is
geometric inertia, since the residue field is algebraically closed.
Thus $\det G=\mathbf F_p^\times$.  Combining this with
\eqref{eq:odd-plane-sl-contained} proves the theorem.

\begin{theorem}\label{thm:plane-odd-all-groups}
Let $p$ be odd and let $U_d$ be the ordinary smooth plane-curve locus
of degree $d$ over any perfect field $k$ of characteristic $p$.  At a
geometric generic point put
$V=H^1_{\mathrm{et}}(C_{\bar\eta},\mathbf F_p)$, of dimension
$g=(d-1)(d-2)/2$.  For every $d\geq3$, both the geometric and the
arithmetic monodromy images on $V$ are
$\mathrm{GL}_g(\mathbf F_p)$, with full determinant character.
For $d=1,2$ both images are trivial on the zero vector space.
For positive rank the representation is irreducible.
\end{theorem}

\begin{proof}
The preceding argument proves the geometric assertion for $d\geq4$
over $\overline{\mathbf F}_p$.  The connected finite etale frame cover
which it identifies is geometrically connected and remains connected
after every algebraically closed field extension.  Since the family
and its Frobenius local system are defined over $\mathbf F_p$, this
gives the same geometric image over the algebraic closure of any
perfect field of characteristic $p$.
The arithmetic image contains that image and is a subgroup of the
same finite general linear group, so it is equal to it.  The cubic
case follows from the same Legendre determinant calculation, since
its rank is one.  Degrees one and two have rank zero.
\end{proof}
